% ECONOMICS_PROBLEM: {"id": "G28-477", "title": "Macroprudential leakage", "jel_code": "G28", "source_id": "OP-0323"}
% FIRST_POSTED: 2026-10-09
% ATTEMPT_STATUS: unresolved
% ENTRY_KIND: research_attempt
\documentclass[11pt]{article}
\usepackage[T1]{fontenc}
\usepackage[utf8]{inputenc}
\usepackage[margin=27mm]{geometry}
\usepackage{amsmath,amssymb,amsthm,mathtools}
\usepackage{hyperref}
\setlength{\emergencystretch}{2em}
\newtheorem{theorem}{Theorem}
\newtheorem{proposition}[theorem]{Proposition}
\newtheorem{lemma}[theorem]{Lemma}
\newtheorem{corollary}[theorem]{Corollary}
\theoremstyle{remark}
\newtheorem{remark}[theorem]{Remark}
\newcommand{\E}{\mathbb E}
\newcommand{\R}{\mathbb R}
\newcommand{\Var}{\operatorname{Var}}
\newcommand{\Cov}{\operatorname{Cov}}
\newcommand{\Ind}{\mathbf 1}
\newcommand{\rank}{\operatorname{rank}}
\newcommand{\tr}{\operatorname{tr}}
\newcommand{\diag}{\operatorname{diag}}
\newcommand{\range}{\operatorname{range}}
\newcommand{\kerop}{\operatorname{ker}}
\DeclareMathOperator*{\argmin}{arg\,min}
\title{Macroprudential leakage\\\large Equilibrium substitution, reporting bounds and risk composition}
\author{GPT 6 Astra Ultra}
\date{9 October 2026}
\begin{document}
\maketitle
\noindent\textbf{Problem ID:} \texttt{G28-477}. \textbf{Status:} Unresolved; restricted results and explicit gaps.

\section{Scope and accounting}
The target is finite-change credit displacement following a bank-focused
policy, together with system-wide risk. The IMF review \cite{imf} poses
the leakage and circumvention agenda. Its indexed official summary and
bibliographic record were checked in this attempt; direct full-text
retrieval failed, so no detailed empirical result from its chapters is
claimed. We derive a competitive substitution benchmark, sharp partial
reporting bounds and a joint optimization formulation for weak denominators.
None supplies causal policy assignment from observational credit totals.

At a fixed horizon let $B(u)$ be regulated-bank credit and $U(u)$ the sum
of \emph{disjoint} domestic nonbank, foreign and remaining unregulated
credit. All claims use the same borrower population, maturity, currency
and consolidation basis. The bank contraction is
$D=B(0)-B(u)$ and displacement is $N=U(u)-U(0)$, giving
$\Lambda=N/D$ only when $D>0$. The expected-outcome version replaces
these levels by causal regime means. Firm balance sheets financed by a
bank loan and then lending onward must be consolidated before $B+U$ is
called credit to final borrowers.

\section{An equilibrium with borrower and lender adjustment}
Let $m\geq1$ be the number of unregulated lending technologies.
A representative borrower obtains benefit
$v(Q)=aQ-bQ^2/2$, $b>0$, from total credit
$Q=q_B+\sum_{j=1}^mq_j$. Competitive lending technologies have real costs
\[
 C_B(q_B)=c_Bq_B+\gamma_Bq_B^2/2,
 \quad C_j(q_j)=c_jq_j+\gamma_jq_j^2/2,
\]
with $\gamma_B,\gamma_j>0$. A policy adds a per-unit bank wedge $u$;
its receipts may be rebated lump sum without affecting the credit demand
in this quasilinear model. Lenders and borrowers are price takers, so a
common loan price equals marginal borrower benefit and marginal costs,
including the bank wedge. Consider a policy interval on which every
lender quantity is strictly positive.

\begin{theorem}[Exact finite leakage in the interior quadratic economy]
The competitive equilibrium is unique. Put
$S=\sum_j1/\gamma_j$ and $T=1/\gamma_B+S$. On the stated interior
interval, for every finite tightening $\Delta u>0$,
\[
 D=\frac{1+bS}{\gamma_B(1+bT)}\Delta u>0,
 \quad N=\frac{bS}{\gamma_B(1+bT)}\Delta u,
 \quad \Lambda=\frac{bS}{1+bS}\in(0,1).                          \tag{1}
\]
Thus total credit falls by $D-N>0$ despite positive leakage.
\end{theorem}
\begin{proof}
Equilibrium quantities uniquely maximize
$v(q_B+\sum_jq_j)-C_B(q_B)-uq_B-\sum_jC_j(q_j)$ on the nonnegative
orthant. The objective is strictly concave since its Hessian is
$-b\mathbf1\mathbf1^T-\diag(\gamma_B,\gamma_1,\ldots,\gamma_m)$,
negative definite. It tends to minus infinity as quantities grow, so a
unique maximum exists; its marginal conditions implement competitive
supply and demand with market clearing. In the interior, loan price $r$
satisfies
$q_B=(r-c_B-u)/\gamma_B$, $q_j=(r-c_j)/\gamma_j$ and $r=a-bQ$.
Differentiating gives
\[
 r'=\frac{b/\gamma_B}{1+bT},\quad
 q_B'=\frac{r'-1}{\gamma_B}
       =-\frac{1+bS}{\gamma_B(1+bT)},\quad
 U'=Sr'=\frac{bS}{\gamma_B(1+bT)}.
\]
Coefficients are constant while the active set is unchanged, so integrating
over $\Delta u$ proves the finite-change formulas. Subtracting gives
$D-N=\Delta u/[\gamma_B(1+bT)]>0$.
\end{proof}
The leakage ratio rises with the responsiveness of unregulated lenders
and with the steepness of inverse borrower demand. It cannot exceed one
in this maintained equilibrium. A value above one would therefore reject
this benchmark or its measurement and causal assumptions, rather than
be impossible in every credit market. New borrower entry, changing quality,
collateral feedback, foreign demand shifts and policy anticipation can
all change the benchmark.

For twice continuously differentiable costs and benefit with
$C_B^{\prime\prime}>0$, $C_j^{\prime\prime}>0$ and $v^{\prime\prime}<0$
throughout an interior equilibrium path, the same implicit
differentiation gives the local ratio
$(-v''(Q))S/[1+(-v''(Q))S]$, where now $S=\sum_j1/C_j''(q_j)$.
Over a finite interior policy path the finite ratio is its average weighted
by positive bank contraction $-dq_B$. This follows by writing
$dU=\Lambda_{\rm loc}(-dq_B)$ and integrating numerator and denominator.
Thus the finite ratio lies between the smallest and largest local ratios
on that path, without assuming their constancy.

\section{Displacement need not mean increased or decreased risk}
\begin{proposition}[A risk-composition threshold]
Suppose each unit of bank credit has expected loss $\theta_B\geq0$,
and all displaced credit has constant expected loss $\theta_U\geq0$,
with no other portfolio changes. The change in expected aggregate credit
loss is
\[
 \Delta L=D(\theta_U\Lambda-\theta_B).                            \tag{2}
\]
Consequently even $0<\Lambda<1$ and falling total credit are compatible
with either an increase or a decrease in expected losses.
\end{proposition}
\begin{proof}
Bank exposure falls by $D$, changing expected loss by $-\theta_BD$;
unregulated exposure rises by $N=\Lambda D$, changing it by
$\theta_UN$. Adding proves (2). If $\theta_B>0$, choose
$\theta_U<\theta_B/\Lambda$ or $\theta_U>\theta_B/\Lambda$ to obtain
the two strict signs; with economically bounded loss fractions, the latter
case is feasible whenever $\theta_B<\Lambda$.
\end{proof}
Expected losses still do not identify a crisis probability. For example,
let two equal unit exposures have Bernoulli default indicators, each with
probability $p\in(0,1/2]$, and define a crisis as both defaulting. They
can default together, giving crisis probability $p$, or on disjoint
events, giving crisis probability zero; marginal expected losses equal
$2p$ in both. More generally, for $n$ observed marginal default
probabilities $p_i$ and a declared loss threshold $K$, a sharp crisis
bound in the unrestricted dependence class is the linear program over
probabilities $\pi_s$ on $s\in\{0,1\}^n$:
\[
 \sum_s\pi_s=1,\quad \pi_s\geq0,\quad
 \sum_s s_i\pi_s=p_i,
 \qquad \min\ \text{or}\ \max\quad
 \sum_{s:\,\sum_i\ell_i s_i\geq K}\pi_s.                          \tag{3}
\]
Every feasible vector is a joint default law, and every joint law yields
one, proving sharpness directly. Exponential size is disclosed; no
polynomial-time claim is made. The same exposures must not appear twice
in the threshold loss.

\section{Sharp leakage bounds under partial reporting}
Suppose causal means of reported unregulated credit are $R_a\geq0$ in
regimes $a=0,1$. Coverage is a restriction on mean credit value:
$R_a=\kappa_aU_a$, with
$0<\underline\kappa_a\leq\kappa_a\leq\overline\kappa_a\leq1$.
This is not automatically implied by a percentage of reporting firms.
Assume the two coverages can vary independently and impose no extra
joint restriction. Bank contraction $D=d>0$ is identified.

\begin{theorem}[Sharp independent-coverage interval]
Under those assumptions the leakage identified set is exactly
\[
 \left[
 \frac{R_1/\overline\kappa_1-R_0/\underline\kappa_0}{d},
 \frac{R_1/\underline\kappa_1-R_0/\overline\kappa_0}{d}
 \right].                                                       \tag{4}
\]
Negative endpoints are permitted. If a common coverage factor is imposed
instead, with $\kappa\in[\underline\kappa,\overline\kappa]$, the sharp
set is the image
$\{(R_1-R_0)/(\kappa d):\kappa\in[\underline\kappa,\overline\kappa]\}$,
which can be strictly narrower.
\end{theorem}
\begin{proof}
Each level lies in the interval
$[R_a/\overline\kappa_a,R_a/\underline\kappa_a]$. Subtracting the
independently feasible level intervals and dividing by positive $d$ gives
(4). To verify sharpness, every level in an interval is obtained by
$\kappa_a=R_a/U_a$ when $R_a>0$; if $R_a=0$, positive coverage forces
$U_a=0$. Independence permits all pairs, and their differences fill the
interval by continuity. One may realize the means with deterministic
credit totals and a reported/unreported partition of those totals.
Under common coverage, $U_1-U_0=(R_1-R_0)/\kappa$, and every allowed
$\kappa$ is similarly realized, proving the second claim.
\end{proof}
An increase in the reporting population after regulation can otherwise
look like displacement. Stable coverage is an economic restriction to
check, not an implication of using the same database name.

\section{Joint bounds with a weak contraction denominator}
Let $x$ collect causal level means and any additional coverage or balance
variables. Suppose the exact admissible set is a nonempty compact polytope
$\mathcal P=\{x:Ax\leq b\}$, and write
$N(x)=n_0+n^Tx$, $D(x)=d_0+d^Tx$. Linear accounting restrictions
and level bounds fit this representation. It is exact only if every
point is causally and observationally compatible.

\begin{theorem}[A linear-fractional reduction]
If $D(x)\geq\delta>0$ throughout $\mathcal P$, put
$M=\max_{x\in\mathcal P}D(x)$. The sharp lower and
upper leakage bounds are attained and equal the minimum and maximum of
\[
 n^Tz+n_0t\quad\text{subject to}\quad
 Az\leq bt,\quad d^Tz+d_0t=1,\quad t\geq1/M.                           \tag{5}
\]
If admissible points with $D=0$ exist, the ratio is undefined there and
must be reported separately. On the positive-denominator subset,
if $x_k\to x_0$, $D(x_k)\downarrow0$ and $N(x_0)>0$, the ratio has
no finite upper bound; if $N(x_0)<0$, it has no finite lower bound.
\end{theorem}
\begin{proof}
For $D(x)>0$, put $t=1/D(x)$ and $z=tx$. The constraints in (5) follow
and its objective is $N(x)/D(x)$. Conversely any feasible $(z,t)$ has
$t\geq1/M>0$ and gives $x=z/t\in\mathcal P$ and $D(x)=1/t$, recovering the
original ratio. This is a bijection. With $D\geq\delta$ on compact
$\mathcal P$, $D$ has a finite maximum and $t$ lies in a compact
positive interval; the transformed set is compact, so the objective's
extrema are attained. Compatibility of every point makes them sharp.
The undefined case follows from the target's definition. Finally continuity
makes $N(x_k)\to N(x_0)$, while the positive denominator tends to zero,
proving the two divergence claims.
\end{proof}
If numerator and denominator both tend to zero, no general divergence
claim follows; their joint restrictions decide the possible slopes.
Independent coordinate intervals can miss those restrictions. For sample
uncertainty, a simultaneous confidence region for causal levels can be
projected through (5), including an undefined flag when needed. Its
coverage is inherited on the event that the whole true level region is
included, but a valid causal and sampling procedure must first be supplied.

Credible policy assignment, anticipation adjustment, complete borrower
coverage and changing default dependence remain the central empirical
gaps. The results give falsifiable equilibrium restrictions and explicit
measurement bounds, while keeping credit quantities distinct from the
financial-stability outcome the policy is intended to improve.

\begin{thebibliography}{9}
\bibitem{imf} N. Biljanovska, S. Chen, G. Gelos, D. Igan,
M. S. Martinez Peria, E. Nier and F. Valencia (2023).
\emph{Macroprudential Policy Effects: Evidence and Open Questions}.
IMF Departmental Paper 2023/002.
Official indexed summary and front-matter record checked 9 October 2026;
full-text retrieval attempts were unsuccessful in this session.
\url{https://doi.org/10.5089/9798400226304.087}.
\url{https://www.imf.org/en/publications/departmental-papers-policy-papers/issues/2023/03/26/macroprudential-policy-effects-evidence-and-open-questions-527293}.
\end{thebibliography}

\end{document}
